\DocumentMetadata{
  pdfversion=2.0,
  pdfstandard=ua-2,
  lang=en-US,
  tagging=on,
  tagging-setup={math/setup=mathml-SE,tabsorder=structure}
}
\documentclass[11pt,letterpaper]{article}
\usepackage[margin=0.68in,includefoot]{geometry}
\usepackage{newcomputermodern}
\usepackage{amsmath,amscd,mathtools}
\providecommand{\square}{\mdlgwhtsquare}
\usepackage[hidelinks]{hyperref}
\hypersetup{
  pdftitle={Numerical Linear Algebra First-Year Exam, January 2018},
  pdfauthor={University of Florida Department of Mathematics},
  pdfsubject={Graduate mathematics examination},
  pdfdisplaydoctitle=true,
  bookmarksopen=true
}
\setcounter{secnumdepth}{0}
\setlength{\parindent}{0pt}
\setlength{\parskip}{0.28em}
\setlength{\emergencystretch}{3em}
\setlength{\leftmargini}{2.15em}
\setlength{\leftmarginii}{2.65em}
\setlength{\leftmarginiii}{3em}
\renewcommand{\labelenumi}{\textbf{\theenumi.}}
\pagestyle{empty}
\raggedbottom
\allowdisplaybreaks
\newenvironment{examproblems}[1]{%
  \begin{enumerate}%
  \setcounter{enumi}{#1}%
  \setlength{\topsep}{0.35em}%
  \setlength{\partopsep}{0pt}%
  \setlength{\itemsep}{0.48em}%
  \setlength{\parsep}{0.18em}%
}{\end{enumerate}}
\newenvironment{examsubparts}{%
  \begin{enumerate}%
  \setlength{\topsep}{0.18em}%
  \setlength{\partopsep}{0pt}%
  \setlength{\itemsep}{0.16em}%
  \setlength{\parsep}{0.1em}%
}{\end{enumerate}}
\newenvironment{examinstructions}{%
  \par\begingroup\itshape
}{%
  \par\endgroup\vspace{0.18em}
}
\makeatletter
\renewcommand\section{\@startsection{section}{1}{0pt}%
  {-1.25ex plus -.25ex minus -.1ex}%
  {0.55ex plus .1ex}%
  {\normalfont\large\bfseries}}
\renewcommand{\@maketitle}{%
  \newpage\null\vskip -1.2em%
  {\footnotesize\bfseries
    \href{https://www.ufl.edu/}{University of Florida}, \href{https://math.ufl.edu/}{Department of Mathematics}\par}%
  \vskip 0.18em%
  \tagstructbegin{tag=Title}%
    {\LARGE\bfseries\href{https://gma.math.ufl.edu/exams/numerical-linear-algebra-first-year/}{Numerical Linear Algebra First-Year Exam}, January 2018\par}%
  \tagstructend%
  \vskip 0.35em%
  \tagmcbegin{artifact}%
    \hrule height 1pt%
  \tagmcend%
  \vskip 0.55em%
}
\makeatother
\title{Numerical Linear Algebra First-Year Exam, January 2018}
\author{}
\date{}
\begin{document}
\maketitle
\thispagestyle{empty}
\begin{examinstructions}
Do 4 (four) problems.
\end{examinstructions}
\begin{examproblems}{0}
\item
Assume \(A \in \mathbb{R}^{m,n}\) with \(m \geq n\), \(\operatorname{rank}(A)=n\), and \(b \in \mathbb{R}^m\).
\begin{examsubparts}
\item[\textnormal{(a)}]
Define the least squares solution to \(Ax=b\).
\item[\textnormal{(b)}]
Derive the normal equations for the least squares problem.
\item[\textnormal{(c)}]
Prove that \(A^T A\) is invertible.
\item[\textnormal{(d)}]
Prove that the unique solution to the least squares problem is \((A^T A)^{-1}A^T b\).
\item[\textnormal{(e)}]
Describe how to solve the least squares problem using the QR decomposition of~\(A\).
\end{examsubparts}
\item
Define a normal matrix and prove that the following are equivalent.
\begin{examsubparts}
\item[\textnormal{(a)}]
\(A\) is normal.
\item[\textnormal{(b)}]
\(\lVert Ax\rVert_2=\lVert A^*x\rVert_2\) for every~\(x\).
\item[\textnormal{(c)}]
\(A\) is unitarily diagonalizable.
\end{examsubparts}
\item
Assume \(A \in \mathbb{R}^{m,m}\).
\begin{examsubparts}
\item[\textnormal{(a)}]
Prove that \(\langle x,y\rangle_A=x^*Ay\) is an inner product on \(\mathbb{R}^m\) if and only if~\(A\) is symmetric and positive definite.
\item[\textnormal{(b)}]
Assume now that~\(A\) is symmetric and positive definite. If \(x_*\) is the solution to \(Ax=b\) and \(\{p_1,\ldots,p_m\}\) is an orthonormal basis for \(\mathbb{R}^m\) with respect to \(\langle\ ,\ \rangle_A\) and \(x_*=\sum c_i p_i\), give a formula for the \(c_i\).
\end{examsubparts}
\item
If \(q_1,\ldots,q_n\) is an orthonormal basis for the subspace \(V \subset \mathbb{C}^m\) with \(m>n\), prove that the orthogonal projector onto~\(V\) is \(QQ^*\), where~\(Q\) is the matrix whose columns are the \(q_j\).
\item
Assume \(A \in \mathbb{C}^{m,m}\).
\begin{examsubparts}
\item[\textnormal{(a)}]
Show that~\(A\) has a Schur decomposition.
\item[\textnormal{(b)}]
If~\(A\) has a collection of~\(m\) linearly independent eigenvectors, show that~\(A\) is diagonalizable.
\end{examsubparts}
\end{examproblems}
\end{document}
